\section{الفصل~\ref{sec:pain}. الألم}

مستشعرات الألم، والسلكان، والبوابة في الحبل الشوكي، ومقبض الصوت الهابط، والتحسّس، والاستيلاء على الانتباه مقيسة مباشرة؛ وصيغتا البوابة والتحسّس نموذجان مبسّطان من الكتاب؛ أما القاعدة التي تختار بها الآلة علوّ صوت الإنذار ففرضية.

{\footnotesize
\setlength{\tabcolsep}{3pt}\renewcommand{\arraystretch}{1.15}
\begin{longtable}{>{\raggedright}p{0.5cm}>{\raggedright}p{4.0cm}>{\raggedright}p{5.6cm}>{\raggedright}p{2.3cm}>{\raggedright\arraybackslash}p{1.7cm}}
\toprule
م & القضية & التجربة وما قيس & المصدر & الحالة \\
\midrule\endfirsthead
\toprule
م & القضية & التجربة وما قيس & المصدر & الحالة \\
\midrule\endhead
1 & عتبة عالية: عتبات التسخين لمستشعرات الألم المختلفة من $41^\circ$ إلى $46\text{--}48^\circ$ & تسجيل ألياف مفردة. جلد القرد: العتبة الوسيطة للمستشعرات البطيئة (\foreignlanguage{english}{C}) $41^\circ$، وللمستشعرات السريعة من النوع الثاني $46^\circ$، ومن النوع الأول فوق $53^\circ$. جلد الإنسان: مستشعرات \foreignlanguage{english}{C} التي تستجيب للضغط أيضًا $40.7^\circ$، والتي لا تستجيب له $48^\circ$. & \cite{Treede1995heat, Weidner1999cfib} & مُقاس \\
2 & مستشعرات الألم تتعوّد أضعف بكثير من مستشعرات اللمس، وتصبح بعد إصابة خفيفة أشد حساسية؛ وضعف الاستجابة بعد تسخين قوي مقيس لدى نوع واحد & حرق جلدي خفيف ($50^\circ$، $100\,\text{s}$): بعد $5\text{--}10$ دقائق تنخفض عتبة الألم لدى البشر وعتبة مستشعرات \foreignlanguage{english}{C} لدى القرد بـ$2\text{--}6^\circ$؛ ولا يوجد هذا الانزياح لدى مستشعرات الجلد الأخرى. ولدى المستشعرات السريعة من النوع الثاني تُكبت الاستجابة بعد التسخين القوي. ومقارنة التعوّد بمستشعرات اللمس تقدير عام، لا تسنده هنا تجربة منفصلة. & \cite{LaMotte1982hyper, Treede1995heat} & مُقاس \\
3 & سلكان: السريع $5\text{--}30\,\text{m/s}$، والبطيء $0.5\text{--}2\,\text{m/s}$؛ وزمن الوصول $t=L/v$~\eqref{eq:paintime} & سرعة التوصيل: مستشعرات الألم السريعة لدى القرد في المتوسط $15$ و$25\,\text{m/s}$ (نوعان)؛ ومستشعرات \foreignlanguage{english}{C} لدى الإنسان $0.8\text{--}1.0\,\text{m/s}$. وعند $L=1\,\text{m}$ يكون ذلك عشرات الملّي ثانية ونحو ثانية. & \cite{Treede1995heat, Weidner1999cfib} & مُقاس \\
4 & يصل الألم مرتين: أولًا السريع الدقيق، ثم الكليل الطويل & زمن رد الفعل لتسخين ساعد الإنسان: من $1100$ إلى $400\ms$ مع ارتفاع الحرارة؛ والتسخين القوي للجلد المشعر يعطي ألمًا مزدوجًا، ولجلد الراحة لا. الألم الأول لا تحمله إلا ألياف أسرع من $6\,\text{m/s}$، ويوافقه بحسب الكمون المستشعرات السريعة من النوع الثاني، وهي غائبة في الراحة. وتقسيم الأدوار (”أين“ و”ما مدى السوء“) تأويل. & \cite{CampbellLaMotte1983first, Treede1995heat} & مُقاس \\
5 & البوابة: عصبون المخرج في الحبل الشوكي يتلقى الألم واللمس، والوسيط المثبِّط يثيره اللمس (الشكل~\ref{fig:gate}) & مخطط البوابة مقترح نظريًا. فئران، وسم جيني وإزالة مجموعات من العصبونات: العصبونات المثيرة الحاملة للسوماتوستاتين لازمة للألم الميكانيكي؛ والعصبونات المثبِّطة الحاملة للداينورفين تمنع دخل مستشعرات اللمس من الوصول إليها، وبدونها تستدعي اللمسة الخفيفة ألمًا. & \cite{MelzackWall1965, Duan2014} & مُقاس \\
6 & اللمس يخفت الألم & الإنسان، نبضات ليزر مؤلمة على اليد: اللمس المتزامن يخفض تقدير الألم والقدرة على تمييز طاقة النبضات؛ ويضعف الأثر مع المسافة بين نقطة اللمس ونقطة الألم داخل القطعة الجلدية الواحدة. & \cite{Mancini2014touch} & مُقاس \\
7 & افتراض نظرية البوابة: الألم القوي يطفئ بنفسه الوسيط المثبِّط، فتنفتح البوابة على مصراعيها & معلَّم في الفصل افتراضًا للنظرية الأصلية. والعمل على الفئران يصف كيف تمنع البوابة اللمس من دخول قناة الألم؛ ولم يُبيَّن فيه إطفاء الوسيط بالألم. & \cite{MelzackWall1965} & فرضية \\
8 & البوابة~\eqref{eq:gate}: العتبة $0.18$ بلا لمس، و$0.86$ مع اللمس، و$0.11$ عند $g=2$؛ وعند $\nu=1$ يخفض اللمس المخرج من $1$ إلى $0.25$، وعند $\nu=2$ لا يؤثر؛ واللمس نفسه لا يمر (الشكل~\ref{fig:gatecurves}) & الحساب بـ\foreignlanguage{english}{\texttt{models/\allowbreak pain\_\allowbreak gate.py}} بأوزان النص يعيد إنتاج كل هذه الأرقام. & \foreignlanguage{english}{\texttt{models/\allowbreak pain\_\allowbreak gate.py}} & نموذج \\
9 & المسارات الهابطة من جذع الدماغ تغيّر كسب البوابة في الاتجاهين & الجرذ، النخاع المستطيل، تسجيل عند تسخين الذيل: بعض العصبونات تفرغ قبل السحب (”تشغيل“)، وبعضها يصمت (”إطفاء“)؛ والتنبيه الضعيف لهذه المنطقة يكبت المنعكس. مراجعة: الدارات نفسها تيسّر نقل الألم وتثبطه، والأفيونيات تعمل عبرها. & \cite{Fields1983rvm, Fields2004} & مُقاس \\
10 & توقع الراحة يخفت الألم عبر القناة الأفيونية & أناس بألم حاد: لدى من خفّض توقعُ الراحة ألمَهم أعاد حجبُ المستقبِلات الأفيونية الألمَ إلى مستوى الآخرين؛ ولدى الآخرين لم يغيّر الحجب الألم. & \cite{Levine1978} & مُقاس \\
11 & يختار الدماغ علوّ صوت الإنذار بحسب جدوى المقاطعة & لا توجد تجربة قيست فيها قاعدة اختيار علو الصوت. & --- & فرضية \\
12 & التحسّس: بعد الإصابة يزداد مخرج البوابة وتنخفض العتبة، جزئيًا بفعل الحبل الشوكي؛ ويدوم حتى بعد توقف المنبّه المؤذي & الجرذ: بعد إصابة الجلد تنخفض عتبة منعكس الثني وتزداد الاستجابة؛ ويبيّن التحليل الكهربائي الفسيولوجي أن جزءًا من الزيادة ينشأ في الحبل الشوكي نفسه. والمنبّه المؤذي يُحدث اضطرابات طويلة في الحساسية تدوم بعد المنبّه نفسه. ولا يذكر الفصل زمن الانحسار الدقيق. & \cite{Woolf1983} & مُقاس \\
13 & التحسّس~\eqref{eq:sensitization} تغذية راجعة موجبة؛ ومع حلقة قوية قد ينغلق الإنذار على نفسه بعد الالتئام & ينتج من النموذج~\eqref{eq:sensitization} (تمرين في آخر الفصل). ولا توجد تجربة تبيّن أن الألم المستمر بعد الالتئام حلقة منغلقة من هذا النوع بالذات. & استنتاج في الفصل & فرضية \\
14 & الألم مكافأة سالبة، والتعلم منه يتبع خطأ الفروق الزمنية & تصوير بالرنين، بشر، سلاسل من نُذر الألم: نشاط الجسم المخطط البطني والقشرة الجزيرية الأمامية يطابق إشارات التعلم بالفروق الزمنية. القرد: بعض عصبونات \nt{DOPA} يُثار بنفخة هواء مزعجة وبنذيرها، وبعضها يُثبَّط. & \cite{Seymour2004, Matsumoto2009} & مُقاس \\
15 & الألم يقاطع العمل الجاري & بشر، منبّه كهربائي مؤلم واختبارات صوتية: الاستجابة لاختبار بعد $100\ms$ من بدء الألم أبطأ بوضوح، وبعد $1.5\,\text{s}$ لا فرق عن الضابط. وتجمع المراجعة تجارب كهذه في نموذج للاستيلاء على الانتباه. & \cite{Crombez1996disrupt, EcclestonCrombez1999} & مُقاس \\
16 & السحب يُغلق في الحبل الشوكي & الجرذ: عصبونات الطبقات العميقة من القرن الخلفي تكرر النمط المكاني لمنعكس السحب لعضلات منفردة؛ ولا يمكن إثارتها بالتنبيه المعاكس الاتجاه من الجزء العنقي، أي أنها وسطاء شوكية. & \cite{Schouenborg1995wr} & مُقاس \\
\bottomrule
\end{longtable}}
